\documentclass{article}
\usepackage[utf8]{inputenc}
\usepackage{amsmath,amssymb}
\usepackage[most]{tcolorbox}
\usepackage{geometry}
\geometry{margin=2.5cm}

\newcommand{\step}[1]{\par\noindent\hangindent=3.5em\hangafter=1%
  \makebox[3.5em][l]{\textbf{#1.}}}
\newcommand{\steptag}[1]{\hfill{\small\textsf{[#1]}}}

\newtcolorbox{proofbox}[1]{
  colback=gray!5, colframe=gray!60,
  fonttitle=\bfseries, title={#1},
  breakable, enhanced,
  left=4pt, right=4pt, top=4pt, bottom=4pt,
  before skip=10pt, after skip=10pt
}

\begin{document}

\begin{proofbox}{Hiddenness depth-Markov: for an exact signed idempotent P...}
\step{1} Hiddenness depth-Markov: for an exact signed idempotent P with delta(P) > 0, nonempty visible set W(P), hidden top vertex v of height H, and any hiddenness dual witness (lambda, alpha, beta) of v with sum\_i beta\_i < kappa = tau/4 (tau = sqrt(delta)), one has for every c > 0: lambda\{f in F\_v : dist\_1(p\_f, conv\{p\_w : w in W\}) > H - c*tau\} > 1 - (1/2 + delta)/c. \steptag{validated} \\[6pt]
\step{1.1} Let C=conv\{p\_w : w in W\} and d\_i=dist\_1(p\_i,C). Because W is nonempty and v is a hidden top vertex of height H, d\_v=H and 0<=d\_i<=H for every row i. Finite-dimensional l1 separation at p\_v gives an affine functional phi with phi(p\_v)=H, phi(x)<=dist\_1(x,C) for all x, and linear part of l\_infty norm at most 1. \steptag{validated} \\[6pt]
\step{1.1.1} By def-height, since W(P) is nonempty and v is a hidden top vertex of height H, for C=conv\{p\_w : w in W\} and d\_i=dist\_1(p\_i,C) one has d\_v=H=max\_i d\_i; therefore 0<=d\_i<=H for every row i. \steptag{validated} \\[4pt]
\step{1.1.2} Finite-dimensional l1 support fact: for any nonempty compact convex C and any point x0 with H0=dist\_1(x0,C), there exists an affine functional phi whose linear part has l\_infty norm at most 1, with phi(x0)=H0 and phi(y)<=0 for every y in C; consequently phi(x)<=dist\_1(x,C) for every x. \steptag{validated} \\[6pt]
\step{1.1.2.1} If H0=0, the zero affine functional has linear part of l\_infty norm 0, satisfies phi(x0)=0, phi(y)=0<=0 on C, and phi(x)=0<=dist\_1(x,C) for all x. \steptag{validated} \\[4pt]
\step{1.1.2.2} If H0>0, compactness of C gives a nearest point y0 in C with ||x0-y0||\_1=H0. \steptag{validated} \\[4pt]
\step{1.1.2.3} For z=x0-y0, the subgradient description of the l1 norm and first-order optimality of y0 for the convex minimization problem min\_\{y in C\} ||x0-y||\_1 give a vector s with ||s||\_infty<=1, s dot z=||z||\_1=H0, and s dot (y-y0)<=0 for every y in C. \steptag{validated} \\[4pt]
\step{1.1.2.4} In the H0>0 case, let y0 be the nearest point from 1.1.2.2 and let s be the support vector supplied by 1.1.2.3; define phi(x)=s dot (x-y0). Then phi(x0)=H0, phi(y)<=0 for y in C, and the linear part has l\_infty norm at most 1. \steptag{validated} \\[6pt]
\step{1.1.2.4.1} Using the validated support-vector node 1.1.2.3 in the H0>0 case, the vector s satisfies ||s||\_infty <= 1, s dot (x0-y0)=H0, and s dot (y-y0) <= 0 for every y in C. Therefore, for phi(x)=s dot (x-y0), phi(x0)=s dot (x0-y0)=H0, phi(y)=s dot (y-y0) <= 0 on C, and the linear part of phi is s, whose l\_infty norm is at most 1. \steptag{archived} \\[4pt]
\step{1.1.2.4.2} By validated node 1.1.2.3, in the H0>0 case the chosen support vector s satisfies ||s||\_infty <= 1, s dot (x0-y0)=H0, and s dot (y-y0) <= 0 for every y in C. Hence for phi(x)=s dot (x-y0), phi(x0)=H0, phi(y)<=0 on C, and the linear part is s with l\_infty norm at most 1. \steptag{validated} \\[4pt]
\step{1.1.2.5} For any x and any y in C, phi(x)=s dot (x-y)+phi(y)<=s dot (x-y)<=||x-y||\_1; taking the infimum over y in C gives phi(x)<=dist\_1(x,C). \steptag{validated} \\[4pt]
\step{1.1.3} Applying that support fact to x0=p\_v and the above C gives the affine phi in 1.1, with phi(p\_v)=H, phi(x)<=dist\_1(x,C) for all x, and linear part of l\_infty norm at most 1. \steptag{validated} \\[6pt]
\step{1.1.3.1} In the context of 1.1, W is nonempty and the row set is finite, so C=conv\{p\_w : w in W\} is a nonempty compact convex subset of finite-dimensional l1 space. By 1.1.1, dist\_1(p\_v,C)=H. Therefore, after the support fact 1.1.2 is validated, universal instantiation of 1.1.2 with x0=p\_v and this C yields an affine phi with phi(p\_v)=H, phi(x)<=dist\_1(x,C) for all x, and linear part of l\_infty norm at most 1. \steptag{validated} \\[4pt]
\step{1.2} With psi=H-phi, every row i satisfies 0<=psi(p\_i), psi(p\_i)>=H-d\_i, and psi(p\_i)<=2+4*delta. \steptag{validated} \\[6pt]
\step{1.2.1} From 1.1, phi(p\_i)<=d\_i<=H for every row i, so psi(p\_i)=H-phi(p\_i) is nonnegative and psi(p\_i)>=H-d\_i. \steptag{validated} \\[4pt]
\step{1.2.2} Because phi(p\_v)=H, psi(p\_i)=phi(p\_v)-phi(p\_i). The linear part of phi has l\_infty norm at most 1, hence psi(p\_i)<=||p\_v-p\_i||\_1. \steptag{validated} \\[4pt]
\step{1.2.3} By the bridge psi(p\_i)<=||p\_v-p\_i||\_1 and def-signed-idempotent row geometry, for an exact signed idempotent with negative mass delta every pair of rows has l1 distance at most 2+4*delta; therefore psi(p\_i)<=2+4*delta. \steptag{validated} \\[6pt]
\step{1.2.3.1} For each row i, node 1.2.2 gives psi(p\_i)<=||p\_v-p\_i||\_1. The row-geometry clause of def-signed-idempotent gives ||p\_v-p\_i||\_1<=2+4*delta for the pair of rows p\_v and p\_i. Therefore psi(p\_i)<=2+4*delta for every row i. \steptag{validated} \\[4pt]
\step{1.3} Using the hiddenness witness balance from lem-hiddenness-dual-witness, applying the linear part of psi gives sum\_f lambda\_f*psi(p\_f)+sum\_i alpha\_i*psi(p\_i)=sum\_i beta\_i*psi(p\_i), hence sum\_f lambda\_f*psi(p\_f)<(1/2+delta)*tau. \steptag{validated} \\[6pt]
\step{1.3.1} By the hypothesis that (lambda,alpha,beta) is a hiddenness dual witness, and by the external lem-hiddenness-dual-witness, lambda\_f>=0 on F\_v with sum\_f lambda\_f=1, alpha\_i,beta\_i>=0, sum\_i beta\_i<kappa=tau/4, and sum\_f lambda\_f*(p\_f-p\_v)+sum\_i alpha\_i*(p\_i-p\_v)=sum\_i beta\_i*(p\_i-p\_v). \steptag{validated} \\[4pt]
\step{1.3.2} Since psi is affine and psi(p\_v)=0, applying the linear part of psi to the witness balance gives sum\_f lambda\_f*psi(p\_f)+sum\_i alpha\_i*psi(p\_i)=sum\_i beta\_i*psi(p\_i). \steptag{validated} \\[4pt]
\step{1.3.3} By 1.2, psi(p\_i)>=0 and psi(p\_i)<=2+4*delta for every row i; with alpha\_i,beta\_i>=0, the equality implies sum\_f lambda\_f*psi(p\_f)<=sum\_i beta\_i*psi(p\_i)<kappa*(2+4*delta). \steptag{validated} \\[4pt]
\step{1.3.4} Because kappa=tau/4, kappa*(2+4*delta)=(1/2+delta)*tau, giving sum\_f lambda\_f*psi(p\_f)<(1/2+delta)*tau. \steptag{validated} \\[4pt]
\step{1.4} Consequently sum\_f lambda\_f*(H-d\_f)<(1/2+delta)*tau. \steptag{validated} \\[6pt]
\step{1.4.1} For each f in F\_v, 1.2 gives psi(p\_f)>=H-d\_f. \steptag{validated} \\[4pt]
\step{1.4.2} The coefficients lambda\_f are nonnegative, so summing the inequalities in 1.4.1 and using 1.3 gives sum\_f lambda\_f*(H-d\_f)<=sum\_f lambda\_f*psi(p\_f)<(1/2+delta)*tau. \steptag{validated} \\[4pt]
\step{1.5} For arbitrary c>0, Markov applied to the nonnegative variables H-d\_f under the probability lambda on F\_v gives lambda\{f in F\_v : d\_f<=H-c*tau\}<(1/2+delta)/c, equivalently lambda\{f in F\_v : d\_f>H-c*tau\}>1-(1/2+delta)/c. \steptag{validated} \\[6pt]
\step{1.5.1} Because delta>0 and tau=sqrt(delta), tau>0; because c>0, c*tau>0. Also, with C=conv\{p\_w : w in W\} and d\_i=dist\_1(p\_i,C), the fact that v is a hidden top vertex of height H gives d\_i<=H for every row i; hence X\_f:=H-d\_f is nonnegative for every f in F\_v. \steptag{validated} \\[6pt]
\step{1.5.1.1} By def-height, since W(P) is nonempty and v is a hidden top vertex of height H, for C=conv\{p\_w : w in W\} and d\_i=dist\_1(p\_i,C) one has H=max\_i d\_i. Therefore d\_f<=H for every row index f, in particular for every f in F\_v, and X\_f=H-d\_f>=0. \steptag{validated} \\[4pt]
\step{1.5.2} Let A\_c=\{f in F\_v : d\_f<=H-c*tau\}. Then A\_c=\{f in F\_v : X\_f>=c*tau\}, so c*tau*lambda(A\_c)<=sum\_\{f in A\_c\} lambda\_f*X\_f<=sum\_f lambda\_f*X\_f. \steptag{validated} \\[4pt]
\step{1.5.3} Using 1.4 and dividing by c*tau gives lambda(A\_c)<(1/2+delta)/c. \steptag{validated} \\[4pt]
\step{1.5.4} Since sum\_f lambda\_f=1 on F\_v, the complement of A\_c has lambda-mass greater than 1-(1/2+delta)/c; this complement is \{f in F\_v : d\_f>H-c*tau\}, which is the asserted set because d\_f=dist\_1(p\_f,conv\{p\_w : w in W\}). \steptag{validated} \\[4pt]
\end{proofbox}

\end{document}
